% Draft method section matching the current implementation.
% Suggested packages: amsmath, amssymb, bm, mathtools, graphicx

\section{Method}
\label{sec:method_revised}

\subsection{Overview and Problem Formulation}
\label{sec:overview_problem}

Given $N$ synchronized views, our input consists of an LDR image and an
event stream for each view,
\begin{equation}
    \mathcal{X}=\{(\bm I_i^{e_i},E_i,\bm K_i)\}_{i=1}^{N},
    \qquad
    E_i=\{(\bm u_m,t_m,p_m)\}_{m=1}^{M_i},
\end{equation}
where $e_i$ is the exposure level, $\bm K_i$ is the camera intrinsic matrix,
and $p_m\in\{-1,+1\}$ is the event polarity. The objective is to recover
multi-view geometry, principally the depth maps $\{D_i\}_{i=1}^{N}$ and their
depth-derived surface normals $\{\bm N_i\}_{i=1}^{N}$. At inference time, the
method requires only the observed LDR images and complete event streams; the
reference exposure and controlled event decomposition described below are
used only during training.

We build on StreamVGGT, which retains the pretrained multi-view RGB
aggregator and geometry heads of VGGT. This backbone provides a strong global
shape prior but loses local evidence when the input is saturated. Our method
adds an event pathway without concatenating raw events with RGB pixels. The
RGB pathway estimates global multi-view structure and absolute depth, whereas
the event pathway supplies exposure-robust temporal evidence and local
differential geometry. Figure~\ref{fig:method_overview} summarizes the complete
architecture.

\subsubsection{Architecture}
\label{sec:architecture}

The RGB images are first processed by the frozen StreamVGGT aggregator,
\begin{equation}
    \{\bm Z_i^{L}\}_{i=1}^{N}
    =\mathcal G_{\mathrm{rgb}}(\{\bm I_i^{e_i}\}_{i=1}^{N}),
\end{equation}
where $\bm Z_i^{L}$ denotes the LDR multi-view tokens. Independently, we
convert every event interval into a polarity-preserving voxel grid with $B$
temporal bins,
\begin{equation}
 V_{i,b}^{p}(\bm u)=
 \sum_m \mathbf 1[\bm u_m=\bm u]\mathbf 1[p_m=p]
 \max(0,1-|b-\bar t_m|),
 \quad
 \bar t_m=(B-1)\frac{t_m-t_0}{t_1-t_0},
\end{equation}
and stack the positive and negative channels as
$\bm V_i\in\mathbb R^{2B\times H\times W}$. Thus, opposite polarities are not
cancelled before encoding. A shared event encoder $\mathcal E$ produces a
dense feature $F_i^E$.

The network then performs two complementary operations. First, a
token-domain HDR adapter uses reliable event evidence to complete the geometry
missing from the LDR tokens,
\begin{equation}
    \bm Z_i^{H}=\bm Z_i^{L}
    +\mathcal A_H(\bm Z_i^{L},\bm T_i^E),
\end{equation}
where $\bm T_i^E$ is obtained by projecting and pooling the selected dense
event feature. A frozen geometry head decodes $\bm Z^H$ into the HDR-like base
depth $D_i^{b}$. Second, a full-resolution event refiner predicts a bounded
log-depth correction $\Delta_i$ from event and coarse-geometry cues,
\begin{equation}
    D_i^{f}=D_i^{b}\exp(\Delta_i),
    \qquad
    \Delta_i=\ell_{\max}\tanh(\mathcal R_D(\cdot)).
\end{equation}
The token path restores global geometry, while the pixel path preserves local
boundaries and surface details that would be removed by patch pooling.

\subsubsection{Observation: Geometry Is a Structured Subset of Events}
\label{sec:event_observation}

An event is triggered by a change of log intensity. Under the standard event
generation model,
\begin{equation}
    \Delta\log I(\bm u,t)\simeq pC_e,
\end{equation}
where $C_e$ is the contrast threshold. For motion-induced changes, the local
brightness variation approximately satisfies
\begin{equation}
    \partial_t\log I(\bm u,t)
    \simeq-\nabla_{\bm u}\log I(\bm u,t)^{\top}\bm v(\bm u,t).
\end{equation}
The image velocity $\bm v$ is coupled to camera motion, depth, surface
orientation, and occlusion boundaries. Consequently, geometry-related events
concentrate around depth discontinuities and normal variations. However, the
complete stream also contains responses caused by reflections, illumination
changes, rendering effects, and sensor noise. Event presence alone is
therefore not equivalent to geometric relevance.

We formalize the complete event stream as a mixture,
\begin{equation}
    \bm V^{\mathrm{full}}
    =\bm V^{\mathrm{geo}}+\bm V^{\mathrm{app}}+\bm V^{\mathrm{noise}},
\end{equation}
where only $\bm V^{\mathrm{full}}$ is available at inference. During training,
controlled event generation additionally provides
$\bm V^{\mathrm{geo}}$, which contains the responses induced by geometric
motion. We do not use this controlled stream as an inference input. Instead,
it defines (i) a geometry-event feature space into which the complete stream
is mapped and (ii) a target for learning spatial event relevance.

This decomposition also clarifies the RGB--event relationship. RGB is not
replaced by events: the RGB backbone supplies global correspondence, camera
context, and absolute-depth structure. Events do not independently regress
the complete scene depth. They first compensate missing LDR token evidence and
then constrain local differential geometry. This division prevents the sparse
event stream from absorbing global scale while allowing it to restore details
that are absent from saturated RGB.

\subsection{Network}
\label{sec:network}

\subsubsection{RGB: Multi-LDR Geometry Consistency}
\label{sec:multi_ldr}

For the same scene, view, and event interval, we sample a severely degraded
observation $I_i^{ev_l}$ and the best available reference exposure
$I_i^{ev_0}$, together with their shared event interval $E_i$. Their
appearance differs, but their underlying geometry is identical. The frozen RGB
aggregator extracts
\begin{equation}
    Z_i^{ev_l}=\mathcal G_{\mathrm{rgb}}(I_i^{ev_l}),
    \quad
    Z_i^{ev_0}=\operatorname{sg}
    [\mathcal G_{\mathrm{rgb}}(I_i^{ev_0})],
\end{equation}
where $\mathcal G_{\mathrm{rgb}}$ is the frozen StreamVGGT geometry encoder and
$\operatorname{sg}$ denotes stop-gradient. In parallel, the event encoder
extracts the complementary dense feature only once,
\begin{equation}
    F_i^E=\mathcal E(E_i).
\end{equation}
The event-conditioned HDR adapter predicts the missing geometry residual
rather than replacing the LDR token,
\begin{align}
    T_i^E
    &=\mathcal P_E(C_i^{\mathrm{fus}}\odot F_i^{E\rightarrow G}),\\
    \widehat Z_i^{H}
    &=Z_i^{ev_l}+\mathcal A_H(Z_i^{ev_l},T_i^E).
\end{align}
In the implementation, $\mathcal A_H$ multiplicatively modulates an LDR-token
context with an event-token context. Its zero-event path is exactly the
identity, so the adapter cannot modify RGB tokens without event evidence.

We supervise the recovered representation with a multi-layer token loss,
\begin{equation}
    \mathcal L_{\mathrm{HDR}}
    =\sum_{l\in\mathcal S}\omega_l
    \rho\!\left(\widehat Z_{i,l}^{H}
    -\operatorname{sg}(Z_{i,l}^{ev_0})\right),
\end{equation}
where $\rho$ is the smooth-$L_1$ penalty and $\mathcal S$ denotes the selected
aggregator layers. Unlike direct feature copying, the residual is conditioned
on the LDR token and selected event evidence. It estimates what is missing
under the current exposure.

The frozen depth head decodes the completed tokens,
\begin{equation}
    D_i^{b}=\mathcal H_D(\widehat Z_i^{H}),
    \qquad
    \bm N_i^{b}=\Pi(D_i^{b},\bm K_i),
\end{equation}
where $\Pi$ is the differentiable depth-to-normal operator. We refer to
$D_i^{b}$ as the HDR-like base or coarse result. Its supervision combines a
depth loss and a geometry-preserving gradient term,
\begin{equation}
    \mathcal L_{\mathrm{base}}
    =\mathcal L_D(D_i^{b},D_i^{\mathrm{gt}})
    +\lambda_g\mathcal L_{\nabla D}(D_i^{b},D_i^{\mathrm{gt}}).
\end{equation}

\subsubsection{Event: Learning Geometry-Aware Event Cues}
\label{sec:event_path}

The event pathway answers two distinct questions: which responses in the
complete stream originate from geometric motion, and how those selected
responses should modify the reconstructed geometry. We address the first with
controlled source attribution and Full-to-Geo transfer, and the second with a
differential-normal guided pixel refiner.

\paragraph{Geometry-First Source Attribution and Full-to-Geo Transfer.}
\label{sec:controlled_event_attribution}

We use the same event encoder for both streams,
\begin{equation}
    \bm F_i^{G}=\mathcal E(\bm V_i^{\mathrm{geo}}),
    \qquad
    \bm F_i^{F}=\mathcal E(\bm V_i^{\mathrm{full}}).
\end{equation}
The controlled geometry stream first trains the encoder to represent local
geometric variation. We then learn a Full-to-Geo aligner $\mathcal A_G$ only
on the complete-stream route,
\begin{equation}
    \widehat{\bm F}_i^{G}=\mathcal A_G(\bm F_i^{F}),
\end{equation}
with the stopped geometry feature as its target,
\begin{equation}
    \mathcal L_{F\rightarrow G}
    =\rho\!\left(\widehat{\bm F}_i^{G}
    -\operatorname{sg}(\bm F_i^{G})\right).
\end{equation}
Separating the two routes is important. Passing
$\bm V^{\mathrm{geo}}$ directly through the aligner would allow the mapping to
learn an identity shortcut, whereas applying the same correction to both
streams would move the teacher space together with the student.

Feature alignment alone may still learn average denoising or noise
compensation. We therefore explicitly predict pixel-level source relevance
from inference-available quantities,
\begin{equation}
    C_i^{\mathrm{fus}}
    =\mathcal Q_F(\bm V_i^{\mathrm{full}},\bm I_i^{L},
                  D_i^{L},\bm N_i^{L})\in[0,1]^{H\times W}.
\end{equation}
The controlled event decomposition supplies the attribution target
\begin{equation}
    C_i^{\mathrm{gt}}(\bm u)=
    \operatorname{clip}\left(
    \frac{\sum_c|V_{i,c}^{\mathrm{geo}}(\bm u)|}
         {\sum_c|V_{i,c}^{\mathrm{full}}(\bm u)|+\epsilon},0,1\right),
\end{equation}
and
\begin{equation}
    \mathcal L_C=\rho(C_i^{\mathrm{fus}}-C_i^{\mathrm{gt}}).
\end{equation}
A second confidence $C_i^{\mathrm{ref}}$ controls the full-resolution detail
path. This separation is necessary because token fusion and pixel refinement
operate in different spaces and need not share the same numerical gate.

Training follows a geometry-first alternating strategy. We first train the
shared event encoder and detail predictor using $\bm V^{\mathrm{geo}}$ while
both confidence maps are fixed. We then alternate between (i) complete-event
attribution, which learns $\mathcal A_G$, $C^{\mathrm{fus}}$, and
$C^{\mathrm{ref}}$, and (ii) Multi-LDR geometry completion, which learns
$\bm I^{L}+\bm V^{\mathrm{geo}}\rightarrow\bm Z^{0}$. At inference,
$\bm V^{\mathrm{geo}}$ is removed and
$\widehat{\bm F}^{G}=\mathcal A_G(\mathcal E(\bm V^{\mathrm{full}}))$ is used.

\paragraph{Differential-Normal Guided Depth Refinement.}
\label{sec:event_normal_depth}

Events primarily reveal changes and boundaries rather than absolute surface
orientation. Predicting a complete normal map from sparse events encourages
the event branch to hallucinate smooth regions. We instead predict the spatial
normal derivative
\begin{equation}
    \widehat{\bm G}_i^{N}
    =\mathcal H_{\nabla N}(\widehat{\bm F}_i^{G})
    \in\mathbb R^{H\times W\times 2\times3},
\end{equation}
corresponding to the horizontal and vertical derivatives of a three-channel
normal. Its target is derived from GT depth,
\begin{equation}
    \bm N_i^{\mathrm{gt}}=\Pi(D_i^{\mathrm{gt}},\bm K_i),
    \qquad
    \bm G_i^{N,\mathrm{gt}}=\nabla\bm N_i^{\mathrm{gt}}.
\end{equation}
The derivative loss emphasizes high-normal-gradient pixels and ignores
unsupported regions,
\begin{equation}
    \mathcal L_{\nabla N}^{E}
    =\frac{\sum_{\bm u}W_i(\bm u)
    \rho(\widehat{\bm G}_i^{N}(\bm u)-
         \bm G_i^{N,\mathrm{gt}}(\bm u))}
    {\sum_{\bm u}W_i(\bm u)+\epsilon},
\end{equation}
where $W_i$ combines valid geometry, genuine event support, recency, and a
larger weight for strong derivative targets. This prevents the numerous
zero-gradient pixels from making an all-zero prediction optimal.

The pixel refiner receives no raw RGB pixels or RGB tokens. Its inputs are the
aligned full-resolution event feature, predicted normal derivative, base log
depth, base normal, and refinement confidence:
\begin{equation}
    \Delta_i=\ell_{\max}\tanh\!\left[
    \mathcal R_D(\widehat{\bm F}_i^{G},
                 \widehat{\bm G}_i^{N},
                 \log D_i^{b},\bm N_i^{b},C_i^{\mathrm{ref}})
    \right].
\end{equation}
The final normal is always recomputed from final depth,
\begin{equation}
    D_i^{f}=D_i^{b}\exp(\Delta_i),
    \qquad
    \bm N_i^{f}=\Pi(D_i^{f},\bm K_i).
\end{equation}
We couple the event derivative to the reconstructed geometry through
\begin{equation}
    \mathcal L_{DN}
    =\frac{\sum_{\bm u}W_i(\bm u)
      \rho(\nabla\bm N_i^{f}(\bm u)-
           \operatorname{sg}(\widehat{\bm G}_i^{N}(\bm u)))}
     {\sum_{\bm u}W_i(\bm u)+\epsilon}.
\end{equation}
An explicit high-frequency residual target further prevents the refiner from
collapsing to a smooth scale correction. Let
\begin{equation}
    R_i^{\mathrm{gt}}=\log D_i^{\mathrm{gt}}-\log D_i^{b},
    \qquad
    R_i^{\mathrm{HF}}=R_i^{\mathrm{gt}}-
    \operatorname{LP}(R_i^{\mathrm{gt}}),
\end{equation}
where $\operatorname{LP}$ is a local low-pass operator. We use
\begin{equation}
    \mathcal L_{\mathrm{HF}}
    =\rho(\Delta_i-R_i^{\mathrm{HF}})
\end{equation}
inside the valid geometry mask.

\paragraph{Adjacent-View Differential-Normal Consistency.}
Independent per-frame supervision does not guarantee that observations of the
same surface produce compatible differential geometry. We therefore exploit
the ordered multi-view input and construct the adjacent pair set
\begin{equation}
    \mathcal P=\{(i,i+1)\mid i=1,\ldots,N-1\}.
\end{equation}
For four input views, this gives the three pairs $(1,2)$, $(2,3)$, and
$(3,4)$. Let $\mathcal W_{i\rightarrow j}$ denote reprojection from view $i$
to view $j$, obtained from the stopped base depth $D_i^b$, camera intrinsics
$K_i,K_j$, and stopped coarse poses. To reduce sensitivity to event sparsity
and small pose errors, we aggregate derivative magnitude and depth within
non-overlapping patches $P$,
\begin{equation}
    g_{i,P}=\operatorname{Avg}_{\bm u\in P}
    \left\|\widehat{\bm G}_i^N(\bm u)\right\|_2,
    \qquad
    \widetilde g_{j\rightarrow i,P}
    =\mathcal W_{j\rightarrow i}(g_{j,P}).
\end{equation}
We compare derivative magnitude rather than its image-axis components because
the magnitude is insensitive to the change of local image axes. A robust
pairwise calibration
\begin{equation}
    a_{ij}=\operatorname{median}_{P\in\Omega_{ij}}
    \frac{\operatorname{sg}(\widetilde g_{j\rightarrow i,P})}
         {\operatorname{sg}(g_{i,P})+\epsilon}
\end{equation}
absorbs the view-dependent derivative scale induced by perspective
projection. The trusted overlap $\Omega_{ij}$ contains only patches with
valid base geometry, event support in both views, in-frame reprojection, and
consistent projected depth. A pair is omitted when either the number or ratio
of overlapping patches is below a threshold. We then define
\begin{equation}
    \mathcal L_{\mathrm{MV}\nabla N}
    =\frac{1}{|\mathcal P_{\mathrm{valid}}|}
    \sum_{(i,j)\in\mathcal P_{\mathrm{valid}}}
    \frac{\sum_{P\in\Omega_{ij}}q_{ij,P}
    \rho\!\left(a_{ij}g_{i,P}-\widetilde g_{j\rightarrow i,P}\right)}
    {\sum_{P\in\Omega_{ij}}q_{ij,P}+\epsilon},
\end{equation}
where
$q_{ij,P}=\sqrt{C_{i,P}^{\mathrm{ref}}
\widetilde C_{j\rightarrow i,P}^{\mathrm{ref}}}$.
The confidence, reprojection depth, and pose are detached in this loss, so the
constraint updates the event derivative branch rather than exploiting a
confidence, depth, or pose shortcut.

\paragraph{Stage-Conditioned Optimization.}
The losses are not optimized as an unstructured parallel sum. They follow the
dependency of the network: the geometry-event representation is learned
first, complete-event attribution is learned second, and the selected event
cues are subsequently used for exposure completion and depth refinement. We
group the objectives as
\begin{align}
    \mathcal L_{\mathrm{geo}}
    &=\mathcal L_{\nabla N}^{E}
      +\lambda_{\mathrm{MV}}\mathcal L_{\mathrm{MV}\nabla N},\\
    \mathcal L_{\mathrm{attr}}
    &=\mathcal L_{F\rightarrow G}+\lambda_C\mathcal L_C,\\
    \mathcal L_{\mathrm{comp}}
    &=\mathcal L_{\mathrm{HDR}}+\lambda_b\mathcal L_{\mathrm{base}},\\
    \mathcal L_{\mathrm{ref}}
    &=\mathcal L_D(D^f,D^{\mathrm{gt}})
      +\lambda_N\mathcal L_N(\bm N^f,\bm N^{\mathrm{gt}})
      +\lambda_{DN}\mathcal L_{DN}
      +\lambda_{\mathrm{HF}}\mathcal L_{\mathrm{HF}}.
\end{align}
The objective is selected according to the alternating training state,
\begin{equation}
    \mathcal L^{(s)}=
    \begin{cases}
      \mathcal L_{\mathrm{geo}}+\mathcal L_{\mathrm{ref}},
      &s=\text{geometry warm-up},\\
      \mathcal L_{\mathrm{attr}},
      &s=\text{complete-event attribution},\\
      \mathcal L_{\mathrm{comp}}+\mathcal L_{\mathrm{ref}}
      +\lambda_{\mathrm{geo}}\mathcal L_{\mathrm{geo}},
      &s=\text{Multi-LDR completion/refinement}.
    \end{cases}
\end{equation}
The confidence maps are fixed during geometry warm-up, while the HDR adapter
and pixel refiner are fixed during source attribution. This ordering prevents
an untrained confidence predictor from suppressing the event branch and
prevents final depth supervision from replacing the controlled attribution
signal.
